Put
\[
J(s_0,s_1)=\Pr\!\left\{s_1(O_{T,1})\le Q_{k_1}(s_1(O_{1,1}),\ldots,s_1(O_{M,1}))\ 
\text{ and }\ s_0(O_{T,0})\le Q_{k_0}(s_0(O_{1,0}),\ldots,s_0(O_{M,0}))\right\}.
\]
By \(\texttt{def:coverage-frontier}\),
\[
\Gamma_{M,\alpha,\boldsymbol\varepsilon}(w)=\max_{s_0,s_1}[(1-\alpha)-J(s_0,s_1)]_+.
\tag{1}
\]
This is exactly the worst deficit over the declared deterministic maps \(s_a:\mathcal O\to[0,K]\). It is not a full-model statement: take fixed covariates, zero predictors, iid nondegenerate Bernoulli variables \(\xi_i\), and \(Y_{iv}(a)=K\xi_i\) for every member and arm. The latent law is bounded and invariant, yet its residual score is nondegenerate conditional on every orbit. Thus the domain in (1) is genuinely narrower than \(\mathcal P_{\mathcal G,K}\), proving the stated full-model disclaimer.

Let \(E_{\mathrm{joint}}\) denote simultaneous armwise inclusion and \(E_{\mathrm{eff}}\) effect-set inclusion. On \(E_{\mathrm{joint}}\),
\[
\tau_T=Y_{0T}(1)-Y_{0T}(0)\in I^{\boldsymbol\varepsilon}_{1,w}-I^{\boldsymbol\varepsilon}_{0,w}=\mathcal C^{\boldsymbol\varepsilon}_w.
\tag{2}
\]
Therefore \(E_{\mathrm{joint}}\subseteq E_{\mathrm{eff}}\), and conditioning on \(\mathcal T\) gives the displayed probability inequality. If both intervals are \(\{0\}\) and both outcomes are \(K\), then the left event fails while \(\tau_T=0\) belongs to the effect set, so equality need not hold.

Assume first that \(k_0,k_1\le M\). For either arm,
\[
s_a(O_{T,a})\le Q_{k_a}(s_a(O_{1,a}),\ldots,s_a(O_{M,a}))
\quad\Longleftrightarrow\quad
\sum_{i=1}^M\mathbf 1\{s_a(O_{i,a})<s_a(O_{T,a})\}<k_a.
\tag{3}
\]
Thus only the total preorder induced by each score map matters. Conversely, every total preorder with ordered nonempty classes \(C_{a,1}\prec\cdots\prec C_{a,d_a}\) has a representative in \([0,K]\): use zero if \(d_a=1\), and otherwise assign \((j-1)K/(d_a-1)\) to class \(j\). Hence the maximum is attained among the \(F_{m_0}F_{m_1}\) preorder pairs.

For every pair \(C\), \(\texttt{lem:coherent-two-finite-rank-frontier-evaluator}\) proves
\[
J(C)=\sum_{\pi\in\Pi_C}\omega_C(\pi)D_{M,k_0,k_1}(\pi),
\tag{4}
\]
while retaining \(R_w\), the coherent target pushforward \(P_T^C\), and within-replicate cross-arm coupling. That lemma and \(\texttt{lem:four-cell-rank-event-dynamic-program}\), \(\texttt{lem:four-cell-multinomial-sum-alternative}\), and \(\texttt{lem:corner-oriented-four-cell-rank-evaluator}\) establish the seven disjoint strata and every displayed exact evaluator. Summing the exact per-key charges is precisely \(\mathsf B\); taking the largest live branch workspace gives \(O(m_0m_1+L+U_{\max}+G+\mathsf W)\). The coherent evaluator proves both time bounds, the factorial-trigger condition, constant workspace for the first four strata, and the edge, face, and rank-one audit identities and resource separations. Every audited selector lookup is normalized and supported on the requested realized arm for each positive-mass assignment; arbitrary total extensions off support have zero integrated weight.

If exactly one rank \(k_b=M+1\), that arm's event is certain. For \(c=1-b\), conditioning (3) on \(O_{T,c}=O_t\) gives a binomial count with success probability
\[
\theta(O_t)=\sum_{O:s_c(O)<s_c(O_t)}q_{c,O}.
\]
Averaging over \(p_{c,O_t}\) proves the displayed one-boundary formula. The class scan and exact shorter-tail evaluation in \(\texttt{lem:one-boundary-orbit-frontier-evaluator}\) give its time and memory bounds. If both ranks equal \(M+1\), both events are certain and (1) equals zero.

It remains to verify the binary-attainment assertion. Take four audited two-member blocks with \(\mathcal G=S_2^4\) and independent fair one-treated--one-control assignment. In each arm choose a block with calibration probabilities
\[
q=(13/20,1/20,1/4,1/20)
\]
and select its unique member observed in the requested arm. This is a normalized, requested-arm-supported kernel on every positive-mass assignment. Give the common target block probabilities \(p=(1/20,1/10,1/4,3/5)\). Use an arm-\(0\) constant score and order the arm-\(1\) blocks as \(1\prec3\prec2\prec4\). At \(M=3\) and \(k_0=k_1=3\), the joint-event failure probability is
\[
\frac14\left(\frac{13}{20}\right)^3+\frac1{10}\left(\frac{18}{20}\right)^3+\frac35\left(\frac{19}{20}\right)^3=\frac{104957}{160000}.
\]
Thus the objective at \(1-\alpha=1/2\) is \(24957/160000\). For binary high-score sets \(H_0,H_1\), conditional independence of the two selector draws gives failure probability
\[
\Phi(H_0,H_1)=p(H_0)\{1-q(H_0)\}^3+p(H_1)\{1-q(H_1)\}^3
-p(H_0\cap H_1)\{1-q(H_0)\}^3\{1-q(H_1)\}^3.
\tag{5}
\]
For an explicit exhaustive certificate, order the sixteen subsets as
\[
\varnothing,1,2,12,3,13,23,123,4,14,24,124,34,134,234,1234.
\]
Substitution of the displayed rational \(p,q\) into (5) gives the following row maxima over \(H_1\):
\[
\left(\frac{20577}{40000},\frac{82651}{160000},\frac{48013}{80000},\frac{20739}{40000},\frac{99183}{160000},\frac{20589}{40000},\frac{25379}{40000},\frac{20579}{40000},\frac{25988367}{40000000},\frac{20723421}{40000000},\frac{3259109}{5000000},\frac{1326351}{2560000},\frac{3259109}{5000000},\frac{164631423}{320000000},\frac{202894331}{320000000},\frac{20577}{40000}\right).
\tag{6}
\]
Each entry is obtained from sixteen evaluations of the explicit polynomial (5), so (6) certifies all \(16^2\) pairs. Its maximum is \(3259109/5000000\), attained at \((H_0,H_1)=(\{2,4\},\{3,4\})\) and its transpose. Therefore every binary objective is at most
\[
\frac{3259109}{5000000}-\frac12=\frac{759109}{5000000},\qquad
\frac{24957}{160000}-\frac{759109}{5000000}=\frac{83189}{20000000}>0.
\]
Binary maps therefore do not attain the frontier.

If target and calibration orbit supports are disjoint in a finite-rank arm, assign score \(K\) to every target-supported orbit and zero to every calibration-supported orbit, and use zero in the other arm. Then \(J=0\), so (1) equals \(1-\alpha\), its largest possible value.

The four cited comparator leaves delimit different optimized objects. \(\texttt{lem:lebrun-rectangular-probability-comparator}\) concerns arbitrary rectangles for discrete count laws; \(\texttt{lem:timans-max-rank-comparator}\) concerns dependence-aware simultaneous family-wise error control; \(\texttt{lem:schlembach-multi-target-allocation-comparator}\) concerns multi-target allocation for hyperrectangle volume; and \(\texttt{lem:structured-volume-dynamic-program-comparator}\) concerns near-optimal structured prediction-set volume. Equations (1)--(4) instead evaluate the present restricted deterministic joint rank-deficit audit with one coherent target. This distinguishes the optimized objects and makes no priority or general-coverage claim.

Finally, \(\texttt{thm:common-selector-lp}\) gives the full-model guarantee under exact balance, and \(\texttt{prop:active-arm-finite-rank-orbit-gap-bound}\) gives the full-model total-variation guarantee under imbalance; neither follows from (1). Setting \(\varepsilon_0=\varepsilon_1=\alpha/2\) gives the equal-split theorem. Since every computational branch evaluates the same probability in (4), none changes the active-arm, pair-audit, empty-library, ordinary-rank, or weighted-orbit semantics.